\documentclass[11pt]{article}
\usepackage[margin=1in]{geometry}
\usepackage{amsmath,amssymb,amsthm}
\usepackage{hyperref}
\usepackage{xurl}
\let\oldus\_
\renewcommand{\_}{\oldus\allowbreak}

\newtheorem{theorem}{Theorem}
\newtheorem{lemma}[theorem]{Lemma}
\theoremstyle{definition}
\newtheorem{definition}[theorem]{Definition}
\theoremstyle{remark}
\newtheorem{remark}[theorem]{Remark}

\newcommand{\V}{\operatorname{V}_p}

\title{A disproof of Conjecture 00000004007\\[2pt]
\large The $p$-variation is $1$-homogeneous, so $\|X+X\| = 2\|X\| > 2^{1/q}\|X\|$}
\date{}

\begin{document}
\sloppy
\maketitle

\section{The conjecture}

Conjecture 00000004007 reads (English):
\begin{quote}
\emph{Definition:} The p-variation is the p-th root of the supremum over partitions of the sum of
p-th powers of the increments. \emph{Conjecture:} The p-variation satisfies the optimal
subadditivity inequality $\|X+Y\| \le (\|X\|^q + \|Y\|^q)^{1/q}$, where $q$ is the conjugate
exponent of $p$, and the inequality cannot be enlarged.
\end{quote}
The conjecture is a conjunction (the inequality holds, and it is optimal). We show that the
inequality itself fails, for every $p\in(1,\infty)$ and also for $p=1$.

\section{Definitions and conventions}

\begin{definition}[Partition]
A partition of $[a,b]$ is a finite sequence $a=t_0<t_1<\dots<t_n=b$ (Lean: \texttt{Partition a b},
fields \texttt{n}, \texttt{t : Fin (n+1) -> R}, \texttt{StrictMono t}, \texttt{t 0 = a},
\texttt{t (last n) = b}).
\end{definition}

\begin{definition}[$p$-variation]
For a path $X:\mathbb{R}\to E$ ($E$ a real normed space; only the values on $[a,b]$ matter) and
$p>0$,
\[
  \V(X;[a,b]) \;=\; \Bigl(\sup_{\pi} \sum_{i=1}^{n} \|X(t_i)-X(t_{i-1})\|^p\Bigr)^{1/p}
  \in [0,\infty],
\]
the supremum over all partitions $\pi$ of $[a,b]$ (Lean: \texttt{incrSum} and \texttt{pVar}).
\end{definition}

Conventions. Values are taken in $[0,\infty]$ (Lean \texttt{ENNReal}), so the supremum always
exists and $\infty^{1/p}=\infty$; our counterexamples have finite values, so this choice is not
used. The conjugate exponent of $p\in(1,\infty)$ is the $q>0$ with $1/p+1/q=1$, i.e.\
$q=p/(p-1)\in(1,\infty)$ (Mathlib's \texttt{Real.HolderConjugate p q}, which requires $p,q>0$ and
$p^{-1}+q^{-1}=1$). For $p=1$ we take $q=\infty$ and read $(a^q+b^q)^{1/q}$ as $\max(a,b)$.
Paths are real-valued in the counterexample; a real-valued counterexample is also one for paths in
$\mathbb{R}^d$ (embed in the first coordinate), and Lemma~\ref{lem:diag} holds in every real normed
space.

\section{Proof}

\begin{lemma}[Homogeneity; \texttt{pVar\_smul}]
For $p>0$, $c\in\mathbb{R}$ and every path $X$: $\V(cX;[a,b]) = |c|\,\V(X;[a,b])$.
\end{lemma}
\begin{proof}
For each partition the increments of $cX$ are $c$ times those of $X$, so each sum is multiplied
by $|c|^p$ (\texttt{incrSum\_smul}). Since $\sup_\pi |c|^p S_\pi = |c|^p\sup_\pi S_\pi$ in
$[0,\infty]$, and $(|c|^p s)^{1/p}=|c|\,s^{1/p}$, the claim follows.
\end{proof}

\begin{lemma}[Failure on the diagonal; \texttt{fails\_at\_diagonal}, \texttt{fails\_at\_diagonal\_one}]\label{lem:diag}
Let $0<\V(X;[a,b])=:v<\infty$. If $1/p+1/q=1$ with $p,q>0$, then
$(v^q+v^q)^{1/q} = 2^{1/q}v < 2v = \V(X+X;[a,b])$. If $p=1$, then $\max(v,v)=v<2v=\mathrm{V}_1(X+X;[a,b])$.
\end{lemma}
\begin{proof}
$X+X=2X$, so the right-hand value is $2v$ by homogeneity. Since $q>1$, $1/q<1$ and $2^{1/q}<2$;
multiply by $0<v<\infty$.
\end{proof}

\begin{lemma}[A non-degenerate path; \texttt{iSup\_incrSum\_id}, \texttt{pVar\_id}]
Let $p\ge1$ and $b-a=1$. For $X(t)=t$, the supremum of the increment sums over partitions of
$[a,b]$ is $1$; in particular $\V(X;[0,1])=1$.
\end{lemma}
\begin{proof}
Each increment $\Delta_i=t_i-t_{i-1}$ lies in $[0,1]$, so $\Delta_i^p\le\Delta_i$ (as $p\ge1$) and
$\sum\Delta_i^p\le\sum\Delta_i=b-a=1$ (telescoping, \texttt{sum\_succ\_sub}). The two-point
partition $a<b$ gives exactly $1$.
\end{proof}

\begin{theorem}[\texttt{conjecture\_00000004007\_false}, \texttt{conjecture\_00000004007\_false'}]
For every $p\in(1,\infty)$ with conjugate exponent $q$, the continuous paths $X=Y=(t\mapsto t)$ on
$[0,1]$ satisfy
\[
  \V(X+Y;[0,1]) = 2 \;>\; 2^{1/q} = \bigl(\V(X;[0,1])^q+\V(Y;[0,1])^q\bigr)^{1/q}.
\]
\end{theorem}

\begin{theorem}[$p=1$; \texttt{conjecture\_00000004007\_false\_one}]
For $p=1$, $q=\infty$: $\mathrm{V}_1(X+Y;[0,1])=2>1=\max(\mathrm{V}_1(X),\mathrm{V}_1(Y))$ for the
same paths.
\end{theorem}

\begin{remark}[Other readings refuted in Lean]
(i) \emph{Raw supremum} (\texttt{raw\_version\_false}): if $\|X\|$ denotes the supremum without the
$p$-th root, it is $p$-homogeneous, the left side at $Y=X$ is $2^p v$ and the right side is
$2^{1/q}v$, and $1/q<1<p$.
(ii) \emph{Concatenation} (\texttt{concat\_version\_false}): if $X+Y$ denotes the concatenation of
$X$ on $[0,1]$ and $Y$ on $[1,2]$ (agreeing at $t=1$), take $Z(t)=t$ on $[0,2]$: both pieces have
$p$-variation $1$ while $\V(Z;[0,2])\ge 2$ by the two-point partition, and $2>2^{1/q}$.
\end{remark}

\section{Scope}

Refuted: the inequality for every $p\in(1,\infty)$ with $q=p/(p-1)$; for $p=1$ with $q=\infty$ read
as the maximum; with or without the $p$-th root; with $X+Y$ the pointwise sum or the concatenation.
Since the inequality is false, the clause ``cannot be enlarged'' is moot. Not addressed:
$0<p<1$ (the conjugate exponent would be negative, and continuous paths of finite $p$-variation are
then constant) and $p=\infty$ (no $p$-th root; $q=1$ and the inequality would reduce to the
triangle inequality for the oscillation, which holds). The counterexample is not degenerate:
$X(t)=t$ is continuous, non-constant and has $p$-variation exactly $1$.

\section{Relation to the earlier submission}

The same mathematical argument ($X=Y$, homogeneity, $2>2^{1/q}$) was merged as PR \#42
(earthking11) and removed in the re-audit of commit 541cf4fb with the reason ``pure powers-of-2
arithmetic; p-variation unformalized''. Its Lean file stated only facts such as
$\neg(2^2\le 2)$ and $2<2^q$ over the natural numbers, and its docstring said ``There is no
real-valued p-variation formalised here''. The present Lean file defines partitions and the
$p$-variation as the $p$-th root of the supremum of the increment sums, proves homogeneity and the
diagonal failure for every path of finite nonzero $p$-variation in every real normed space,
computes the $p$-variation of $t\mapsto t$ exactly, and states the conjecture's inequality with
Mathlib's conjugate-exponent predicate for every real $p>1$ (not only integer instances).

\section{Lean formalization}

File \texttt{lean/Conjecture4007/Basic.lean} (Lean 4.33.1, Mathlib), namespace \texttt{Tlmc4007}.
Main theorems: \texttt{conjecture\_00000004007\_false} (for all $p,q$ with
\texttt{p.HolderConjugate q}), \texttt{conjecture\_00000004007\_false'} (for all $p>1$ with
$q=p/(p-1)$), \texttt{conjecture\_00000004007\_false\_one}, \texttt{raw\_version\_false},
\texttt{concat\_version\_false}; supporting \texttt{pVar\_smul}, \texttt{fails\_at\_diagonal},
\texttt{pVar\_id}. \texttt{lake build} succeeds; \texttt{\#print axioms} of each main theorem gives
\texttt{[propext, Classical.choice, Quot.sound]}. No \texttt{sorry}, no \texttt{native\_decide}, no
added axioms.

\end{document}
